\documentclass[11pt]{article}
\usepackage[margin=1in]{geometry}
\usepackage{amsmath,amssymb}
\usepackage{graphicx}
\usepackage{booktabs}
\usepackage{xcolor}
\usepackage[hidelinks]{hyperref}

% Figures live in each scenario directory (each owns its plots); pull them from there.
\graphicspath{%
  {../air_bgo_100kBq/figures/}%
  {../air_bgo_1MBq/figures/}%
  {../air_bgo_10MBq/figures/}%
  {../water_bgo_1MBq/figures/}%
  {../../nema/nema_la_air_bgo/figures/}%
  {../../nema/nema_la_water_bgo/figures/}%
}

\newcommand{\sr}{s_{r}}
\newcommand{\sens}{\text{sens}}
\newcommand{\code}[1]{\texttt{#1}}

\title{Reconstructing PTCRYSP Monte-Carlo Data:\\
       Normalization, Randoms, Attenuation, Scatter on a Uniform Sphere,\\
       and Reconstruction Noise on a NEMA Contrast Phantom}
\author{RecoCrysp --- MC reconstruction studies (Part III)}
\date{\today}

\begin{document}
\maketitle

\begin{abstract}
We reconstruct Monte-Carlo listmode coincidence data from the PTCRYSP simulation
of the CRYSP scanner, using the \code{RecoCrysp} listmode Poisson MLEM. A uniform
sphere phantom is the test object throughout. Because every coincidence carries a
ground-truth class flag (true / scatter / random) and a per-gamma scatter count,
we can build and \emph{validate} each correction --- normalization, randoms,
attenuation, scatter --- against the truth. The recurring lesson is that a uniform,
centred sphere is spatially degenerate with any smooth additive background, so the
randoms and scatter corrections, though demonstrably correct, are nearly invisible
on this phantom; attenuation, being multiplicative and radially structured, is the
one effect the sphere displays dramatically. We then move to a NEMA NU-2
image-quality contrast phantom, where the spheres \emph{are} displayed, and use it
to track down the reconstruction noise. The headline result is that the dominant
background mottle is neither the data statistics nor the voxel size but the
Monte-Carlo sampling noise of the \emph{sensitivity image} $A^{\!\top}(1)$:
under-sampled it injects a $\sim\!30\%$ background coefficient of variation,
amplified by the MLEM denominator, that falls as $1/\sqrt{n_{\text{sens}}}$ and is
removed by computing the sensitivity from a large LOR sample. Normalization, in
other words, is itself a sizeable Monte-Carlo job. Finally, the contrast phantom in
water --- attenuation plus $20.5\%$ scatter and $3.2\%$ randoms together --- is where
the additive corrections become visible: the backgrounds wash the sphere contrast and
the correction recovers part of it, the demonstration the degenerate uniform sphere
could not give.
\end{abstract}

\section{Setup and method}

\paragraph{Data.} PTCRYSP writes one row per coincidence (\code{lors\_det.h5}):
the two detector endpoints, the emission origin, the energy, the detector element
indices, and --- crucially --- a truth flag $\in\{0,1,2\}$ (true/scatter/random)
and the per-gamma Compton-scatter counts $n_{\text{scat}1},n_{\text{scat}2}$. Single
versus multiple scatter is $n_{\text{scat}1}+n_{\text{scat}2}=1$ versus $\ge 2$. A
reader (\code{RecoExamples}) parses the schema directly into endpoint matrices and
boolean class masks; \code{tools/lors\_summary.py} prints the class breakdown of any
dataset and is the standing first step on each new file.

\paragraph{Scanner and phantom.} The CRYSP scanner is a continuous, monolithic BGO
cylinder: inner radius $387$\,mm, wall $37$\,mm, half-length $512$\,mm (axial FOV
$1024$\,mm), read out as $48\,\varphi\times 20\,z = 960$ elements. The phantom is a
uniform $R=80$\,mm sphere of $^{18}$F at the origin, simulated in vacuum (air) or
water.

\paragraph{Model.} The expected counts on LOR $i$ are
\begin{equation}
  \bar y_i \;=\; n_i\,a_i\,(A\,x)_i \;+\; s_i \;+\; r_i,
  \label{eq:model}
\end{equation}
with $n$ the normalization, $a$ the attenuation survival, $A$ the (Joseph) forward
projector, $x$ the activity image, and $s,r$ the scatter and randoms backgrounds.
In the listmode Poisson MLEM each event is its own LOR; the multiplicative factor
$\text{mult}=n\,a$ enters the forward model and the additive $\text{contamination}=s+r$
enters as background:
\begin{equation}
  x \leftarrow \frac{x}{\text{sens}}\;A^{\!\top}\!\!\left(\frac{\text{mult}\cdot\text{counts}}{\text{mult}\,(Ax)+\text{contam}}\right),
  \qquad
  \text{sens}=A^{\!\top}(\text{mult}).
  \label{eq:mlem}
\end{equation}
The sensitivity is source-free: at the MLEM fixed point the source cancels in the
data ratio, so $\text{sens}$ must enumerate the detectable LORs once. We sample them
on the scanner surface and put the bare back-projection on the per-event scale with
$\text{scale}=n_{\text{events}}/n_{\text{sens}}$. All kernels run unchanged on CPU
or Apple Metal (the latter $\sim\!40\times$ faster per iteration here).

All reconstructions in this report use one locked recipe, matched to the NEMA study
of Section~\ref{sec:nema}: a $2.5$\,mm grid; the sensitivity sampled at
$n_{\text{sens}}=5\times10^{8}$ LORs (the count needed to push its Monte-Carlo
sampling noise below the floor --- the central finding of Section~\ref{sec:nema},
the reason an apparently trivial $A^{\!\top}(1)$ is a large MC job); and a
regularized update in place of plain MLEM (Part~I). The noise control is chosen per
scenario along two axes --- does the sensitivity get small (attenuation), and is there
fine contrast to preserve:
\begin{itemize}
  \item \emph{Vacuum, contrast} (NEMA in air) --- edge-preserving \textbf{Huber} (OSL):
        bounds the noise while keeping the small-sphere edges.
  \item \emph{Water, no contrast} (uniform sphere) --- \textbf{De Pierro} quadratic
        smoothness (\code{penalized\_mlem}), because OSL Huber \emph{diverges} under
        attenuation: there $\text{sens}=A^{\!\top}(a)$ is small at the centre, so the
        one-step-late denominator $\text{sens}+\nabla R$ collapses and a hot voxel runs
        away. De Pierro is monotone and cannot; with no contrast to lose, its lack of
        edge preservation costs nothing.
  \item \emph{Water, contrast} (NEMA in water, Section~\ref{sec:nemawater}) ---
        \textbf{MLEM + Gaussian post-filter}. Quadratic smoothness would erase the hot
        spheres, and OSL Huber would diverge; plain MLEM ($x^{+}=t/\text{sens}$, no
        division by $\text{sens}+\nabla R$) is stable under attenuation, and a post-filter
        controls the noise while preserving contrast (post-filter $\approx$ Huber on the
        NEMA frontier, Section~\ref{sec:nema}).
\end{itemize}

\section{Normalization (air, $n=1$)}

For a continuous monolithic crystal there is no per-crystal efficiency variation,
so the library's analytic geometric sensitivity $A^{\!\top}(1)$ \emph{is} the
normalization: $n\equiv 1$, provided the analytic scanner geometry matches the
simulated one. Reconstructing the vacuum sphere with $n=1$ and no contamination
yields a quiet, near-flat interior (Fig.~\ref{fig:norm}): the per-voxel coefficient
of variation is $\approx 0.10$, on top of a radial tilt (centre low, rim high).

That residual texture is \emph{not} Poisson noise --- an early read of this run, on a
coarse grid and an under-sampled sensitivity, put the interior CoV at $\approx 0.29$
and called it Poisson, which the NEMA study (Section~\ref{sec:nema}) later showed was
wrong: it was the Monte-Carlo sampling noise of $A^{\!\top}(1)$ itself, sampled at
only $2\times10^{7}$ LORs. At the locked $5\times10^{8}$ the CoV drops to $\approx
0.10$, exactly the NEMA behaviour, and the interior is genuinely flat in noise.

The radial tilt is a separate, real effect. At the locked recipe it is larger than
the $\pm 8\%$ once quoted: the converged profile runs from a centre at $\approx 0.74$
of the interior mean up to $\approx 1.06$ near the rim --- a $\sim\!25\%$ centre
depression. The increase is \emph{not} the regularizer: an unregularized MLEM scan
(\code{diagnostics/scan\_niter.jl}) converges to the same centre/peak ratio
($0.72$ at $30$ iterations, versus $0.69$ for Huber), and it is not under-stopping
either (the ratio settles by $\sim\!20$ iterations). It is the finer $2.5$\,mm grid
resolving a tilt the old $4$\,mm voxels averaged away: at matched iterations the only
change from the $\pm 8\%$ run is the voxel size. The sign and origin are unchanged ---
feeding the pipeline ideal, efficiency-free events reproduces the same ratio, so the
tilt is a LOR-sampling \emph{measure} mismatch (surface-endpoint-pair sampling versus
the emission-detection measure of the real LORs), not a detector efficiency needing an
MC normalization scan. We carry it as a fixed baseline; which measure is the correct
one to normalize with remains the open question (the diagnostics live in
\code{air\_bgo\_100kBq/diagnostics/}).

\begin{figure}[t]
  \centering
  \includegraphics[width=\linewidth]{air_bgo_100kBq}
  \caption{Normalization on the vacuum sphere (locked recipe: $n=1$, Huber,
  $5\times10^{8}$ sensitivity LORs, $2.5$\,mm). Left: the true uniform sphere
  (central slice). Middle: the reconstruction (same slice), flat in noise
  (interior CoV $\approx 0.10$). Right: the radial profile, with the $\sim\!25\%$
  centre-low measure tilt that the fine grid resolves.}
  \label{fig:norm}
\end{figure}

\section{Randoms (air, activity series)}

Randoms scale with \emph{activity}, not event count. Three vacuum datasets at the
same statistics but $100$\,kBq, $1$\,MBq, and $10$\,MBq (over $1000$, $100$, and
$10$\,s) give random fractions of $0.03\%$, $0.94\%$, and $8.72\%$ --- linear in
activity, as expected ($r_{ij}=2\tau\,S_i S_j/T$). The randoms model is the same
singles-free, class-agnostic estimator used for scatter and for the NEMA randoms
study (Section~\ref{sec:nema}): the flagged random LORs (truth $=2$) are histogrammed
in sinogram coordinates, smoothed, and read back per event as the local random
fraction, calibrated so its sum equals the flagged random total (it matches exactly,
e.g.\ $544{,}451$ at $1$\,MBq). No singles file is needed.

The three-way comparison (trues-only gold / uncorrected / corrected) runs cleanly
at the locked recipe, but even at $8.7\%$ the corrected and uncorrected
reconstructions of the uniform sphere differ by under a percent in the interior
(slice mean $|$uncorr$-$corr$|\approx 0.7\%$): the randoms form a smooth, near-uniform
pedestal that is degenerate with the sphere and is divided out by the interior-mean
normalization. The only clean signature is the outside-sphere halo, which the
correction pulls partway back (pedestal gold $0\to$ uncorr $0.0018\to$ corr $0.0014$
at $10$\,MBq). The machinery is validated but the effect stays invisible on this
phantom; a visible randoms demonstration needs a contrast phantom, which is why the
NEMA study (Section~\ref{sec:nema}) carries it. (Drivers and figures under
\code{air\_bgo\_1MBq/} and \code{air\_bgo\_10MBq/}.)

\section{Attenuation (water, trues + AC)}

Water introduces real attenuation. For a uniform sphere the survival factor of a
LOR is analytic --- it depends only on the chord $\ell$ the line cuts through the
sphere, not on where the annihilation occurred:
\begin{equation}
  a_i = \exp(-\mu\,\ell_i), \qquad \mu_{\text{water}}(511\,\text{keV}) = 0.096~\text{cm}^{-1}.
\end{equation}
We reconstruct the \emph{true} coincidences (scatter removed by the flag, so the
narrow-beam $\mu$ is exactly the value to correct with) twice: without correction
($\text{mult}=1$) and with ($\text{mult}=a$, $\text{sens}=A^{\!\top}(a)$).

Without correction the image cups severely (Fig.~\ref{fig:att}): the reconstructed
centre/edge ratio is $\approx 0.4$, above the line-integral prediction
$\exp(-\mu\,2R)=0.215$ (the excess is finite-resolution smearing), with the rim
$\sim\!3\times$ the centre. The correction inverts it to a flat interior ---
centre/edge $\approx 1.1$, the inside recovered to within a few percent of uniform
(normalized $0.87$--$1.03$), the small residual being the measure tilt seen in
vacuum. This is a cleaner recovery than the original coarse-grid MLEM run (which
reached only $0.80$): the locked recipe, with the De Pierro quadratic prior holding
the centre stable where $A^{\!\top}(a)$ is weakest, flattens the sphere fully.
Attenuation, unlike the additive backgrounds, is multiplicative and radially
structured, so the sphere displays it fully.

\begin{figure}[t]
  \centering
  \includegraphics[width=\linewidth]{water_bgo_1MBq_att_only}
  \caption{Attenuation correction on the water sphere (true coincidences). Left:
  no correction --- the textbook radial cupping. Middle: attenuation-corrected.
  Right: radial profiles; the correction removes the cupping, leaving the vacuum
  measure tilt.}
  \label{fig:att}
\end{figure}

\section{Scatter (water)}

Scatter is the dominant contamination here: $17.2\%$ of coincidences ($15.8\%$
single, $1.4\%$ multiple). It enters as the additive $s$ in Eq.~\eqref{eq:model}.
Having the flags, we build a \emph{smooth} scatter model from the flagged scatter
events (as a single-scatter simulation would), rather than subtracting individual
events. Each LOR is reduced to sinogram coordinates --- radial offset $\sr$, axial
midpoint $z_m$, obliquity $\Delta z$ --- and the per-event background is the smoothed
local scatter fraction $s_i = \tilde S_{b(i)}/\tilde P_{b(i)}$ (smoothed scatter
counts over smoothed prompt counts in the bin), which sums to the scatter total by
construction.

\paragraph{Inspect the model before trusting it.} The scatter fraction in $\sr$ has
a genuinely sharp physical edge at $\sr=R$: inside the sphere it is $\approx 0.13$,
and beyond $R$ --- where no true LOR can land --- it must approach $1$
(Fig.~\ref{fig:sino}). Our first Gaussian smoothing blurred this step (the halo
fraction came out $0.56$ instead of $\sim\!1$), and an auto-ranged $\sr$ axis wasted
resolution on the empty far field. Capping the ranges for finer bins and smoothing
only $z_m,\Delta z$ (leaving $\sr$ sharp) restores a halo fraction of $0.997$. The
lesson: smoothing must respect the real edges of the distribution.

\begin{figure}[t]
  \centering
  \includegraphics[width=\linewidth]{water_bgo_1MBq_sinogram}
  \caption{Scatter-model sinograms. The prompt counts $P$ are confined to the
  sphere footprint ($\sr\le R$); the scatter $S$ spreads well beyond. Bottom-right:
  the scatter fraction versus $\sr$ has a sharp step at $\sr=R$ (dotted line) that
  the edge-preserving smoothing keeps intact.}
  \label{fig:sino}
\end{figure}

\paragraph{Result.} The scatter model calibrates exactly to the flagged total
($3{,}964{,}440$ events) and its action is physical but, on this phantom at the locked
recipe, almost entirely confined to where scatter dominates (Fig.~\ref{fig:scat}, and
the zoned diagnosis \code{att\_scatter\_zones.py}): the whole sphere interior --- core,
mid, both edges --- moves by under $0.1\%$, while the pure-scatter far halo is removed
by $-9\%$ and the corners by $\sim\!-84\%$ ($0.032\to0.005$). The radial-profile
distance to the gold standard is unchanged ($0.143$ corrected versus $0.143$
uncorrected): inside the sphere the scatter is degenerate with the source and the
interior-mean normalization divides out whatever uniform pedestal survives, so the
three profiles overlap. The correction is doing exactly the right thing --- subtracting
scatter where it lives --- but the uniform sphere cannot \emph{display} it. (At the
finer grid the correction is even more invisible than the original coarse run, which
had seen a small edge ring and a profile-distance improvement; the stronger smoothing
and the better-resolved degeneracy erase it.)

\begin{figure}[t]
  \centering
  \includegraphics[width=\linewidth]{water_bgo_1MBq_att_scatter}
  \caption{Scatter correction on the water sphere, all attenuation-corrected. Top:
  gold (trues only), uncorrected (trues+scatter), scatter-corrected slices. Bottom:
  radial profiles; the uncorrected$-$corrected difference map (robust scale, dotted
  circle $=R$) showing the edge ring and far-halo removal; and the difference from
  gold versus radius.}
  \label{fig:scat}
\end{figure}

\section{A contrast phantom: NEMA image quality and the source of the noise}
\label{sec:nema}

The recurring limitation above is the phantom: a uniform sphere cannot display a
correction that is degenerate with it. We therefore move to a NEMA NU-2
image-quality phantom --- a $R=100$\,mm body cylinder containing six hot spheres
(diameters $10,13,17,22,28,37$\,mm) on a $57.2$\,mm ring at a $4{:}1$
sphere-to-background activity ratio --- reconstructed in vacuum on the same BGO
scanner (\code{nema\_la\_air\_bgo}, low activity, $2.9\%$ randoms; the phantom,
its sphere and background ROIs, and the truth map are built by
\code{src/nema\_phantom.jl}). The figure of merit is the contrast recovery
coefficient
\begin{equation}
  \text{CRC} = \frac{\bar x_{\text{sphere}}/\bar x_{\text{bg}} - 1}{R_{\text{true}} - 1},
  \qquad R_{\text{true}} = 4,
  \label{eq:crc}
\end{equation}
($100\%$ = full contrast recovered), paired with the background coefficient of
variation $\text{CoV}=\sigma/\mu$ over a uniform region, the standard NEMA
noise measure.

\paragraph{The background is mottled, and it is not the obvious culprits.} An
early-stopped MLEM of this phantom shows the spheres but on a heavily speckled
background --- a per-voxel CoV of $\approx0.30$, far above the clinical NEMA range
($\sim\!5$--$15\%$). We chased it through every knob in turn, each as a saved scan
script, and the first two suspects came up empty:
\begin{itemize}
  \item \emph{Counts / Poisson statistics --- not the driver.} Subsampling the true
        coincidences from $N/8$ to $N$ at fixed regularization leaves the
        background CoV essentially flat ($0.32\!\to\!0.30$ over an $8\times$ range),
        where Poisson noise alone would change it by $\sqrt8\approx2.8\times$. More
        events do not quiet this background (\code{scan\_counts.jl}).
  \item \emph{Voxel size --- not the driver.} Scanning the voxel from $1.5$ to
        $4$\,mm at a fixed post-filter leaves the CoV flat ($0.32\!\to\!0.29$, finer
        slightly worse), so the mottle is not a grid-sampling effect
        (\code{scan\_voxel.jl}).
\end{itemize}
What the noise \emph{does} track is the iteration count: the CoV grows roughly
linearly with iteration while the mean converges, the signature of an ill-posed
problem whose late iterations fit high-frequency detail (\code{scan\_niter.jl}).
Regularization bounds that growth, and on the small-sphere CRC-versus-CoV frontier
a quadratic smoothness prior, an edge-preserving Huber prior
(\code{penalized\_mlem} / \code{osl\_mlem}, Part~I), and a plain Gaussian
post-filter on an early-stopped MLEM all sit on nearly the same curve
(Fig.~\ref{fig:nema-method}): post-filter $\approx$ Huber $>$ quadratic $\approx$
unfiltered MLEM. Edge preservation gives no clear gain here because the small
spheres are low-contrast --- there are no strong edges for the Huber threshold to
protect. Regularization controls the noise but does not explain it; the CRC/CoV
trade is also degenerate (it rewards over-smoothing), so it cannot be free-optimized.

\begin{figure}[t]
  \centering
  \includegraphics[width=0.78\linewidth]{nema_la_air_bgo_method_compare}
  \caption{Small-sphere contrast recovery versus background noise for four noise
  controls at matched settings. Post-filtered MLEM and the edge-preserving Huber
  prior trace nearly the same frontier; the quadratic prior and unfiltered MLEM are
  comparable. For these low-contrast spheres edge preservation gives no clear
  advantage --- the noise is ill-conditioning that any of the four can bound, so
  the method choice is close to a wash and we simply lock one for the studies.}
  \label{fig:nema-method}
\end{figure}

\paragraph{The source: the sensitivity image is a Monte-Carlo estimate.} The knob
that was held fixed in all of the above is the one that mattered. The sensitivity
$\sens=A^{\!\top}(1)$ is not analytic here: it is back-projected from a finite
sample of $n_{\text{sens}}$ geometric LORs drawn on the scanner surface, so it
carries its own $\sim\!1/\sqrt{n_{\text{sens}}}$ Monte-Carlo sampling noise --- and
$\sens$ sits in the \emph{denominator} of every MLEM update, where that noise is
amplified. Varying $n_{\text{sens}}$ with everything else fixed confirms it
(Fig.~\ref{fig:nema-nsens}, \code{scan\_nsens.jl}): the background CoV falls as
$0.30\,(20\text{M})\to0.18\,(50\text{M})\to0.12\,(100\text{M})\to0.10\,(200\text{M})
\to0.08\,(500\text{M})$, tracking the $1/\sqrt{n_{\text{sens}}}$ line and staying a
fixed $\approx5\times$ above the sensitivity image's own CoV --- the conditioning
amplification. So the noise \emph{is} statistics, but the statistics of the
sensitivity estimate (held at $20$\,M LORs in every earlier scan), not of the data.
The fix is to compute $\sens$ from a large LOR sample; it is a one-time cost. At
$500$\,M the background settles near a $0.03$--$0.06$ floor and the six spheres
stand out cleanly with CRC rising from $43\%$ on the smallest to $106\%$ on the
largest (Fig.~\ref{fig:nema-spheres}); the low-$n_{\text{sens}}$ CRC was itself
inflated by the sensitivity mottle. The lesson generalizes: for a continuous
scanner whose normalization \emph{is} the geometric sensitivity, normalization is a
sizeable Monte-Carlo job, not a cheap analytic prefactor.

\begin{figure}[t]
  \centering
  \includegraphics[width=0.62\linewidth]{nema_la_air_bgo_nsens_scan}
  \caption{The background noise is the sensitivity-image sampling. Background CoV
  and the sensitivity-image CoV versus the number of sensitivity LORs, everything
  else fixed. Both fall as $1/\sqrt{n_{\text{sens}}}$; the background tracks the
  sensitivity CoV at a constant $\approx5\times$ amplification (the MLEM
  ill-conditioning). The $20$\,M sampling used in the count/voxel scans was the
  hidden noise source.}
  \label{fig:nema-nsens}
\end{figure}

\begin{figure}[t]
  \centering
  \includegraphics[width=\linewidth]{nema_la_air_bgo_spheres_500M}
  \caption{NEMA reconstruction with the sensitivity sampled at $500$\,M LORs (MLEM,
  $6$\,mm post-filter). Left: the central slice, all six hot spheres clean on a
  quiet background (CoV $0.08$). Right: contrast recovery versus sphere diameter,
  $43\%$ (smallest) to $106\%$ (largest).}
  \label{fig:nema-spheres}
\end{figure}

\section{Full physics on a contrast phantom: NEMA in water}
\label{sec:nemawater}

The water NEMA phantom (\code{nema\_la\_water\_bgo}) closes the arc: a contrast object
\emph{and} all three corrections at once --- attenuation, $20.5\%$ scatter, $3.2\%$
randoms --- where, unlike the uniform sphere, the corrections are no longer degenerate
with the source and can be \emph{seen}. The body is a uniform water cylinder
($R=100$\,mm, $|z|\le90$\,mm), so attenuation is again exact and analytic, now via the
\emph{cylinder} chord (the line's radial interval $x^2+y^2\le R^2$ intersected with the
axial slab; \code{cylinder\_chord}) rather than the sphere chord. Reconstruction is
MLEM\,+\,$8$\,mm post-filter: quadratic smoothness erased the hot spheres outright and
OSL Huber diverges under attenuation, so the stable, contrast-preserving choice is
post-filtered MLEM (the per-scenario reasoning is in the method note above). We reconstruct three ways,
all attenuation-corrected: \emph{gold} (true coincidences), \emph{uncorrected} (all
prompts, no scatter/randoms model), and \emph{corrected} (contamination $=$ smoothed
sinogram scatter $+$ randoms).

The corrections move the numbers (Fig.~\ref{fig:nemawater}). The scatter and randoms
fill the cold background and wash the contrast: the $37$\,mm sphere's CRC falls from the
gold $74\%$ to $48\%$, and the background mean rises from $0.467$ to $0.514$. The
correction pulls both back toward gold --- CRC $48\to50\%$, background $0.514\to0.501$ ---
and the three CRC-versus-size curves separate cleanly (gold well above; corrected just
above uncorrected at every size). This is the visible, quantitative correction the
uniform sphere could not provide. The recovery is real but \emph{partial}: most of the
gold-to-uncorrected gap remains, as the smoothed additive model under-subtracts in the
cold background. Tuning the scatter-model parametrization for stronger cold-region
subtraction is the natural refinement; the machinery and the geometry are validated.

\begin{figure}[t]
  \centering
  \includegraphics[width=\linewidth]{nema_la_water_bgo}
  \caption{Full physics on the water NEMA phantom (MLEM\,+\,$8$\,mm post-filter,
  $500$\,M sensitivity LORs). Top: central slices --- gold (trues, AC), uncorrected
  (all prompts, AC only), corrected (AC\,+\,scatter\,+\,randoms); all six hot spheres
  reconstruct on the water body. Bottom: contrast recovery versus sphere diameter for
  the three. The curves separate --- scatter\,+\,randoms wash the contrast (gold$\to$
  uncorrected) and the correction recovers part of it (uncorrected$\to$corrected) ---
  the demonstration the degenerate uniform sphere could not give.}
  \label{fig:nemawater}
\end{figure}

\section{Synthesis}

\begin{table}[t]
  \centering
  \begin{tabular}{@{}llll@{}}
    \toprule
    Effect & Character & On the uniform sphere & Validated \\
    \midrule
    Normalization & geometric, $n=1$    & flat in noise (CoV $0.10$), $\sim\!25\%$ measure tilt & yes \\
    Randoms       & smooth additive     & degenerate, invisible            & yes \\
    Attenuation   & multiplicative      & strong cupping, fully corrected  & yes \\
    Scatter       & smooth additive     & degenerate; far-field residual   & yes \\
    \bottomrule
  \end{tabular}
  \caption{The four corrections on the uniform sphere. The additive backgrounds
  (randoms, scatter) are validated but nearly invisible; the multiplicative,
  radially structured attenuation is the one the sphere displays.}
  \label{tab:summary}
\end{table}

Across all four studies the through-line is the same (Table~\ref{tab:summary}): a
uniform, centred sphere is spatially degenerate with smooth additive backgrounds,
and normalizing radial profiles to the interior mean removes whatever residual
survives. The corrections are correct --- the estimates calibrate exactly to the
flagged totals, and the structured, physical action is visible in absolute,
un-normalized terms --- but the phantom cannot \emph{show} them. Attenuation is the
exception precisely because it is multiplicative and radially structured.

The NEMA contrast phantom is where the sphere's degeneracy lifts. In vacuum
(Section~\ref{sec:nema}) it gives a quantitative CRC and --- the main reconstruction
finding of this report --- exposes the background noise as the Monte-Carlo sampling of
the sensitivity image, fixable by sampling $\sens$ from a large LOR set rather than by
more data or a finer grid. In water (Section~\ref{sec:nemawater}) it finally
\emph{shows} the additive corrections: with attenuation, $20.5\%$ scatter and $3.2\%$
randoms all present, the scatter\,+\,randoms wash the contrast and the correction
recovers part of it, the CRC curves separating where the uniform sphere kept them on
top of one another. The natural continuations are to push the cold-background
subtraction of the scatter model (the recovery is currently partial), and to run the
high-activity randoms twin ($22.9\%$) for a randoms-only control. Two smaller
side-studies remain self-contained: relaxing the energy window on the water dataset
(which would raise the $0.13$ interior scatter fraction), and settling which LOR
measure is the correct one for the $\sim\!25\%$ normalization tilt.

\end{document}
